% hwfig user manual.  Compile with pdflatex (twice, for the contents).
% Needs hwfig.sty on the search path; uses ../examples/ and snippets/.
\documentclass[11pt,a4paper]{article}
\usepackage[T1]{fontenc}
\usepackage{lmodern}
\usepackage[margin=2.3cm]{geometry}
\usepackage{booktabs,array,listings,xcolor}
\usepackage{hwfig}
\usepackage[hidelinks]{hyperref}

\lstset{basicstyle=\ttfamily\small,columns=fullflexible,keepspaces,
  frame=single,rulecolor=\color{black!25},xleftmargin=4pt,xrightmargin=4pt,
  framexleftmargin=4pt,framexrightmargin=4pt,upquote=true}
\newcommand\code[1]{\texttt{\detokenize{#1}}}
\newcommand\cs[1]{\texttt{\textbackslash#1}}
\newcommand\snippet[1]{%
  \lstinputlisting{snippets/doc-#1.tex}\par\medskip
  \begin{center}\input{snippets/doc-#1.tex}\end{center}}
\setlength\parindent{0pt}\setlength\parskip{0.6em}

% one strip per palette: the eight role swatches and the three wire colours
\newcommand\palstrip[1]{%
  \begin{tikzpicture}[hw,hwcanvas]\hwpalette{#1}
    \node[anchor=west,font=\ttfamily\small,text=hwtext] at (-0.2,0) {#1};
    \foreach \r [count=\i] in {io,comp,quant,const,op,mem,ctrl,group}{%
      \node[draw=hw\r D,fill=hw\r B,minimum width=0.95cm,minimum height=0.5cm,
            font=\tiny,text=hw\r T] at ({1.6+\i*1.05},0) {\r};}
    \draw[hwwire,line width=1pt] (11.3,0.12) -- ++(0.9,0);
    \draw[hwcwire,line width=1pt,dashed] (11.3,0) -- ++(0.9,0);
    \draw[hwaccent,line width=1.4pt] (11.3,-0.12) -- ++(0.9,0);
  \end{tikzpicture}\par}

\input{../examples/catalogue.inc}

\title{\textbf{hwfig}\\[2pt]\large TikZ styles for hardware and dataflow figures}
\author{Pelayo Leguina\\\small\url{https://github.com/pleguina/hwfig}}
\date{Version 1.2, 15 November 2026}

\begin{document}
\maketitle

\begin{center}
\includegraphics[width=\linewidth]{../examples/l1_trigger.pdf}
\end{center}

\tableofcontents

\section{What it is for}

Block diagrams of FPGA firmware, HLS kernels and trigger systems tend to be
drawn from scratch for every paper, and they rarely look alike. hwfig gives
you a fixed set of TikZ styles for the things that appear in these figures
again and again: compute kernels, quantisers, compile-time constants,
memories, FIFOs, processors, AI~Engine tiles, routers, clock domains and so
on.

Every style takes its colours from a \emph{role} (input/output, compute,
quantisation, constant, operator, memory, control, frame) instead of from a
fixed colour. That makes three things easy:

\begin{itemize}
  \item switching a whole paper between print-friendly greyscale and a
        colour-blind-safe palette with one line;
  \item changing the overall look (rounded or square boxes, solid or tinted
        fills, serif or sans) without touching the figures;
  \item keeping every figure in a document consistent, because nobody picks
        colours by hand.
\end{itemize}

hwfig only needs TikZ and xcolor. It works with pdfLaTeX, XeLaTeX and LuaLaTeX.

\section{Installation}

If hwfig is in your TeX distribution, there is nothing to do. Otherwise copy
\code{hwfig.sty} next to your document (on Overleaf, upload it to the
project), or put it in your personal tree so every document finds it:

\begin{lstlisting}
mkdir -p "$(kpsewhich -var-value TEXMFHOME)/tex/latex/hwfig"
cp hwfig.sty "$(kpsewhich -var-value TEXMFHOME)/tex/latex/hwfig/"
\end{lstlisting}

hwfig adds one pgf layer, \code{hwcanvas}, below the usual
\code{background} layer. If your document calls \cs{pgfsetlayers} itself,
keep \code{hwcanvas,background,main} in the list.

\section{A first figure}

Load the package and start each picture with the \code{hw} style. It sets
the font size, the line width and the arrow tips.

\begin{lstlisting}
\usepackage{hwfig}            % greyscale, classic look
\end{lstlisting}

\snippet{first}

Run-time data uses \code{runtime}, values fixed at compile time (weights,
scale factors, configuration) use the dashed \code{ctime}, and the thick
\code{hot} wire marks the path you want the reader to follow.

\section{Palettes}

Choose a palette as a package option, \code{\usepackage[okabe]{hwfig}}, or
switch anywhere with \cs{hwpalette}\code{{okabe}}. The switch is local to
the current group, so one figure can use a different palette from the rest
of the document.

\medskip
\palstrip{grey}\palstrip{okabe}\palstrip{tol}\palstrip{bright}
\palstrip{ibm}\palstrip{nord}\palstrip{slate}\palstrip{cern}
\palstrip{ngt}\palstrip{ngtblue}\palstrip{ngtdark}\palstrip{ngtnight}
\medskip

\begin{description}
  \item[\code{grey}] The default. Safe for print and photocopies.
  \item[\code{okabe}, \code{tol}, \code{bright}, \code{ibm}] Colour-blind-safe
    palettes (Okabe--Ito, Paul Tol's muted and bright sets, IBM's design
    palette). \code{okabe} has the strongest contrast and is a good default
    for colour papers.
  \item[\code{nord}, \code{slate}] Softer tones that suit slides.
  \item[\code{cern}] CERN blue with an orange accent.
  \item[\code{ngt}, \code{ngtblue}] The colours of CERN's Next Generation
    Triggers project: the five logo colours, and the navy, teal and amber of
    its title slides.
  \item[\code{ngtdark}, \code{ngtnight}] The same two palettes on a dark
    background. Add \code{hwcanvas} to the picture options,
    \code{\begin{tikzpicture}[hw,hwcanvas]}, so the background is painted.
\end{description}

The three lines at the end of each strip are the normal wire, the constant
wire and the accent used by \code{hot}. A misspelt palette name stops the
compilation with an error rather than silently falling back.

\subsection{Your own palette}

A palette sets a draw colour and a light fill for each of the eight roles,
plus the three wire colours. Keep the fills light (at most about 25\,\%) so
black text stays readable.

\snippet{palette}

A dark palette also calls \cs{hwink}\code{{<text>}{<paper>}} to set the text
and background colours, and mixes its fills into the background colour
rather than into white.

\section{The look: shapes, fills, fonts and lines}

The look is independent of the palette and, like it, local to the current
group. \cs{hwsetup} sets several things at once, and the same keys work as
package options:

\begin{lstlisting}
\usepackage[ngt,shape=round,font=sans]{hwfig}
\hwsetup{palette=ngtdark,shape=pill,fill=solid,line=bold,shadow}
\end{lstlisting}

\begin{center}
\begin{tabular}{>{\ttfamily}l p{0.62\linewidth}}
\toprule
\normalfont command & values \\
\midrule
\textbackslash hwshape & \code{classic} (default: square kernels, rounded
  ports), \code{soft}, \code{round}, \code{pill} (pill-shaped ports),
  \code{sharp} \\
\textbackslash hwfill  & \code{tint} (default), \code{solid} (saturated boxes
  with light text), \code{outline} \\
\textbackslash hwfont  & \code{serif}, \code{sans}, \code{mono}, or any font
  commands \\
\textbackslash hwline  & \code{thin}, \code{normal}, \code{bold}, or a length \\
\textbackslash hwshadow & a soft drop shadow under boxes; \cs{hwshadow}\code{[false]}
  turns it off \\
\bottomrule
\end{tabular}
\end{center}

\newcommand\looksample[1]{%
  \begin{tikzpicture}[hw]\hwsetup{#1}
    \node[anchor=west,font=\ttfamily\scriptsize,text=black!70] at (-0.7,0.75) {\detokenize{#1}};
    \node[runtime] (x) at (0,0) {$x$};
    \node[hlsbox,minimum width=1.5cm] (k) at (2,0) {MAC};
    \node[qbox] (q) at (4.1,0) {sat8};
    \node[mem,minimum width=1.3cm] (m) at (6.2,0) {buffer};
    \node[ctrl,minimum width=1.2cm] (f) at (8.2,0) {FSM};
    \draw[arr] (x)--(k); \draw[hot] (k)--(q); \draw[arr] (q)--(m); \draw[ctlarr] (f)--(m);
  \end{tikzpicture}\par\smallskip}
\begin{center}\hwpalette{okabe}
\looksample{shape=classic}
\looksample{shape=round}
\looksample{shape=pill,font=sans}
\looksample{shape=round,fill=solid}
\looksample{shape=sharp,fill=outline,line=bold}
\end{center}

Use \code{solid} with a colour palette; with \code{grey} it gives black
boxes. For sans-serif maths as well as text, load the \code{sansmath}
package and use \cs{hwfont}\code{{\sffamily\sansmath}}.

\clearpage
\section{Node and wire styles}

\subsection{Nodes}

\begin{center}
\small
\begin{tabular}{>{\ttfamily}l p{0.6\linewidth}}
\toprule
\normalfont style & use \\
\midrule
runtime, ctime & run-time inputs and outputs; compile-time constants \\
hlsbox, qbox, op & compute kernel; quantiser or requantiser; arithmetic operator \\
mux, demux & multiplexer and demultiplexer \\
mem, bram, uram, fifo & on-chip memory and streams \\
reg, ff & register with clock wedge \\
lut, dsp & FPGA primitives \\
extmem, cpu, host & DDR or HBM; processor; host computer \\
xbar, router, pe, aie & interconnect, NoC router, processing element, AI Engine tile \\
chip, port, cdc & package with pins, pin marker, clock-domain crossing \\
term, decision & flow-chart end point and branch \\
state, accstate & FSM state and accepting state \\
buf, net, neuron & driver or buffer, network cloud, neural-network unit \\
stack & add to any box to show several copies of it \\
dot, badge, note & wire junction, numbered call-out, free annotation \\
group, clkdomain & background frames (use with \code{fit}) \\
plregion, psregion & frames for the programmable logic and processor sides \\
grouplabel, lbl, lblbox & frame titles and wire labels \\
\bottomrule
\end{tabular}
\end{center}

\subsection{Wires}

\code{arr} is the normal data wire and \code{darr} the dashed wire from a
constant. \code{hot} highlights a path, \code{bus} is a wide datapath (add
\code{pic[pos=.5]{buswidth=64}} for the width), \code{biarr} goes both ways,
\code{ctlarr} carries control or handshakes, and \code{link} is a serial or
optical link. \code{brace} and \code{brace'} draw a dimension brace along a
segment.

The catalogue below shows every style in the \code{grey} and \code{ngt}
palettes.

\begin{center}
\def\pal{grey}\hwpalette{grey}\resizebox{\linewidth}{!}{\catalogue}
\end{center}
\begin{center}
\def\pal{ngt}\hwpalette{ngt}\resizebox{\linewidth}{!}{\catalogue}
\end{center}

\section{Parametric pictures}

Some structures are drawn by a single \cs{pic}. Options go in the pic call
and only affect that pic; \cs{hwset}\code{{...}} changes the defaults for
the rest of the current scope instead. The parts inside a pic get names
starting with the pic's own name, so you can connect wires to them.

\begin{center}
\small
\begin{tabular}{>{\ttfamily}l p{0.3\linewidth} p{0.36\linewidth}}
\toprule
\normalfont pic & options & named parts \\
\midrule
aiearray & rows, cols, dx, dy, tilelabel & tiles \code{n-t<r>-<c>}, shims \code{n-s<c>}, frame \code{n-box} \\
systolic & rows, cols & \code{n-pe<r>-<c>} \\
addertree & depth & inputs \code{n-in<k>}, adders \code{n-a<l>-<k>}, \code{n-out} \\
pipeline & labels & stages \code{n-st<k>}, registers \code{n-r<k>} \\
noc & rows, cols, nocpe & routers \code{n-n<r>-<c>}, PEs \code{n-pe<r>-<c>} \\
fabric & columns, fabricrows & column tops \code{n-c<k>} \\
mlp & layers, lsep, nsep & units \code{n-l<layer>-<k>} \\
legend & legend & --- \\
\bottomrule
\end{tabular}
\end{center}

\snippet{pics}

\section{Checking a layout}

While you draft, load the package with the \code{check} option, or put
\cs{hwsetup}\code{{check}} at the start of one picture. hwfig then looks for
the mistakes that are easy to miss on screen:

\begin{itemize}
  \item two boxes that overlap;
  \item a wire label that sits on a box;
  \item a wire that runs through a box;
  \item a wire too short for its arrow head, or one that turns right at the
        arrow head.
\end{itemize}

Each problem gets a warning in the log and a red mark in the PDF, and every
picture ends with a line such as \code{hwfig check: 0 problem(s)}. Add
\code{nocheck} to a node that overlaps another on purpose, and remove the
option for the final version. Here the check finds an overlap, a wire
through two boxes and an arrow that is too short:

\snippet{check}

The check cannot judge everything: wires crossing each other, labels on
frame borders and general crowding still need a look at the result.

\section{Drawing conventions}

These are the conventions the examples follow. They are not enforced, but
they keep figures readable.

\begin{itemize}
  \item Data flows from left to right. Keep the main path on one row and
        put parallel branches above and below it.
  \item Place compile-time constants directly above or below the block they
        configure, joined with \code{darr}.
  \item Follow every block that changes precision with a \code{qbox}.
  \item Place nodes at explicit coordinates on a grid (about 2.4\,cm between
        columns); this is easier to keep tidy than relative positioning.
  \item Route wires with straight and right-angled segments, never through
        boxes, and leave at least 0.4\,cm before each arrow head.
  \item Put labels only on stretches of wire with free space around them.
  \item Leave at least 0.6\,cm between neighbouring frames, and keep figures
        within about 16\,cm for a two-column paper.
\end{itemize}

\section{Examples}

The \code{examples} folder has complete documents: a quantised GNN layer, a
Versal SoC, the Level-1 trigger on the first page, a track-finder FPGA, an
NGT-style trigger chain, and the catalogues of styles, palettes and looks.

\section{License and history}

hwfig is distributed under the \LaTeX\ Project Public License, version 1.3c
or later. The example documents may also be used under CC~BY-SA~4.0.
Changes between versions are listed in \code{CHANGELOG.md}, and bug reports
and suggestions are welcome on GitHub.

\end{document}
